\section*{Abstract}
		Die T0-Theorie untersucht, wie weit sich kosmische Phänomene auf eine einzige universelle Konstante $\xi = \frac{4}{3} \times 10^{-4}$ zurückführen lassen. Dieses Dokument präsentiert die fundamentalen Beziehungen zwischen der Gravitationskonstante, der kosmischen Mikrowellenhintergrundstrahlung (CMB), dem Casimir-Effekt und kosmischen Strukturen im Rahmen eines statischen, ewig existierenden Universums. Alle Herleitungen erfolgen in natürlichen Einheiten ($\hbar = c = k_B = 1$) und stützen sich auf die korpuseigene Zeit-Energie-Reziprozität $T\cdot E = 1$ (Dok.~306).

	\medskip\noindent\textbf{Hinweis zur Fassung.} Die deutsche Fassung dieses Dokuments ist umfangreicher als die englische. Bei früheren Überarbeitungen wurden in der englischen Fassung Teile gekürzt oder anders gefasst. Bei Abweichungen gilt die umfangreichere deutsche Fassung.


	\medskip\noindent\textbf{Statusmarker.} Aussagen mit \textbf{[S]} sind \emph{spekulativ}: ein Hinweis oder eine Vermutung, die im Korpus noch nicht hergeleitet oder numerisch geprüft ist. Dieses Dokument behandelt fast ausschließlich den kosmischen Sektor (statisches Universum, CMB, Rotverschiebung, Dunkle Energie). Der heutige Korpus klammert diesen Sektor bewusst aus, weil auch das Standardmodell ihn nicht herleitet (P39); die kosmologischen Folgerungen dieses Dokuments sind daher als \textbf{[S]} zu lesen.

	
	
	\section{Einführung: Die universelle $\xi$-Konstante}
	
\subsection{Grundlagen der T0-Theorie}

\begin{important}
	Die T0-Theorie basiert auf der universellen dimensionslosen Konstante $\xi = \frac{4}{3} \times 10^{-4}$, auf die sie Phänomene vom subatomaren bis zum kosmischen Bereich zurückführt.
\end{important}

Die T0-Theorie baut auf einer einzigen fundamentalen Konstante auf. Diese Konstante bildet die Grundlage für alle physikalischen Berechnungen und Vorhersagen der Theorie:

\begin{equation}
	\xi = \frac{4}{3} \times 10^{-4} = 1.333333... \times 10^{-4}
\end{equation}

Diese dimensionslose Konstante verbindet Quanten- und Gravitationsphänomene und ist als Grundlage einer einheitlichen Beschreibung der fundamentalen Wechselwirkungen gedacht.

\begin{tcolorbox}[colback=yellow!10!white,colframe=yellow!50!black,title=Hinweis zur Herleitung]
	Für die detaillierte Herleitung und physikalische Begründung dieser fundamentalen Konstante siehe das Dokument ``Parameterherleitung'' (verfügbar unter: ).
\end{tcolorbox}

	\subsection{Zeit-Energie-Dualität als Fundament}
	
	\begin{revolutionary}[These]
		\textbf{[S]} Eine klassische Urknall-Singularität mit unendlicher Dichte ist physikalisch nicht realisiert; an ihre Stelle tritt ein winziger, aber endlicher Kern mit minimaler Längenskala $L_0$ (aus $\xi$). Die Begründung liefert die korpuseigene Zeit-Energie-Reziprozität $T\cdot E = 1$ (Dok.~306).
	\end{revolutionary}
	
	In natürlichen Einheiten ($\hbar = c = k_B = 1$) folgt aus der Zeit-Masse-Dualität $T \cdot m = 1$ mit $E = m$ die exakte Reziprozität:
	
	\begin{equation}
		T \cdot E = 1
	\end{equation}
	
	Daraus ergeben sich die kosmologischen Konsequenzen \textbf{[S]}:
	\begin{itemize}
		\item Aus $T \cdot E = 1$ folgt $E \to \infty \Leftrightarrow T \to 0$
		\item Die Zeit ist nach unten durch $T_0 = \xi\,t_P$ beschränkt; damit ist $E$ nach oben beschränkt, $E \le 1/T_0$
		\item Eine unendliche Energiedichte ist strukturell ausgeschlossen: es gibt keinen singulären Anfang
		\item Der Torus ist kompakt und randlos ($\partial T^4 = \emptyset$); die Frage nach einem Anfang bei $t = 0$ setzt einen Rand voraus, den er nicht besitzt
		\item Das Universum ist statisch, ohne expandierenden Raum
	\end{itemize}
	
	Heisenbergs Unschärferelation $\Delta E \times \Delta t \geq \frac{1}{2}$ trägt dieses Argument nicht: Bei endlichem $\Delta t$ liefert sie nur eine endliche untere Schranke für $\Delta E$, und ihr $\Delta t$ ist die Änderungszeit einer Observablen, kein kosmisches Alter (Dok.~306).

	\section{Kosmische Mikrowellenhintergrundstrahlung (CMB)}
	
	\subsection{CMB ohne Urknall: $\xi$-Feld-Mechanismen}
	
	\begin{revolutionary}[These]
		\textbf{[S]} Wenn die Zeit-Energie-Dualität einen Urknall ausschließt, wie oben gefolgert, muss die CMB einen anderen Ursprung haben als die z=1100-Entkopplung der Standardkosmologie.
	\end{revolutionary}
	
	Die T0-Theorie deutet die CMB \textbf{[S]} als Strahlung von $\xi$-Feld-Quantenfluktuationen:
	
	\begin{equation}
		\frac{T_{\text{CMB}}}{E_\xi} = \frac{16}{9} \xi^2
	\end{equation}
	
	Mit $E_\xi = \frac{1}{\xi} = \frac{3}{4} \times 10^4$ (natürliche Einheiten) und $\xi = \frac{4}{3} \times 10^{-4}$ ergibt sich:
	
	\begin{equation}
		T_{\text{CMB}} = \frac{16}{9} \xi^2 \times E_\xi = \frac{16}{9} \times 1{,}78 \times 10^{-8} \times 7500 = 2{,}37 \times 10^{-4}
	\end{equation}
	
	\textbf{Umrechnung in SI-Einheiten:}
	\begin{equation}
		T_{\text{CMB}} = 2{,}751 \text{ K}
	\end{equation}
	
	Der Wert $2{,}370\times10^{-4}$ entspricht mit $E_\xi$ in eV $2{,}751$~K; gemessen sind $2{,}725$~K. Der Ausdruck liegt damit etwa 1\,\% neben dem beobachteten Wert (vgl.\ die Tabelle der dimensionslosen Verhältnisse: $3{,}13$ gegenüber $3{,}16 \times 10^{-8}$). Diese Übereinstimmung ist einheitenabhängig: Die Relation ergibt eine reine Zahl, die erst in der Energieeinheit eV mit dem Messwert übereinstimmt; sie ist keine Herleitung, die Frage ist offen.
	
	\subsection{CMB-Energiedichte und $\xi$-Längenskala}
	
	Die CMB-Energiedichte in natürlichen Einheiten beträgt:
	\begin{equation}
		\rho_{\text{CMB}} = \frac{\pi^2}{15}T_{\text{CMB}}^4 \approx 2{,}0 \times 10^{-15}\ \text{eV}^4 \quad \text{(natürliche Einheiten, Dimension } [E^4] \text{)}
	\end{equation}
	
	Diese Energiedichte definiert eine charakteristische $\xi$-Längenskala:
	\begin{equation}
		L_\xi = \left(\frac{\xi}{\rho_{\text{CMB}}}\right)^{1/4}
	\end{equation}
	
	\begin{formula}
		Fundamentale Beziehung der CMB-Energiedichte:
		\begin{equation}
			\rho_{\text{CMB}} = \frac{\xi}{L_\xi^4} = \frac{\frac{4}{3} \times 10^{-4}}{(L_\xi)^4}
		\end{equation}
	\end{formula}
	
	\section{Casimir-Effekt und $\xi$-Feld-Verbindung}
	
	\subsection{Casimir-CMB-Verhältnis}
	
	\begin{experiment}
		Das Verhältnis zwischen Casimir-Energiedichte und CMB-Energiedichte bei der charakteristischen $\xi$-Längenskala $L_\xi = 10^{-4}$ m ist eine Identität, weil $L_\xi$ selbst aus der gemessenen CMB-Energiedichte und $\xi$ bestimmt wird; es ist berechnet, kein Messwert.
	\end{experiment}
	
	Die Casimir-Energiedichte bei Plattenabstand $d = L_\xi$ beträgt:
	\begin{equation}
		|\rho_{\text{Casimir}}| = \frac{\pi^2}{720 \times L_\xi^4} \quad \text{(natürliche Einheiten)}
	\end{equation}

	Die Formel $\pi^2\hbar c/(240\,d^4)$ beschreibt den Casimir-Druck zwischen idealen Platten, die Energiedichte ist $\pi^2\hbar c/(720\,d^4)$.
	
	Das Verhältnis ergibt:
	\begin{equation}
		\frac{|\rho_{\text{Casimir}}|}{\rho_{\text{CMB}}} = \frac{\pi^2}{720 \xi} = \frac{\pi^2 \times 10^4}{960} \approx 102{,}8
	\end{equation}

	\textbf{Direkte Berechnung mit SI-Werten:}
	Mit $L_\xi = 10^{-4}$ m ergibt die direkte Berechnung:
	\begin{align}
		|\rho_{\text{Casimir}}| &= \frac{\hbar c \pi^2}{720 \times (10^{-4})^4} = 4{,}33 \times 10^{-12} \text{ J/m}^3 \\
		\rho_{\text{CMB}} &= 4{,}17 \times 10^{-14} \text{ J/m}^3 \\
		\text{Verhältnis} &= \frac{4{,}33 \times 10^{-12}}{4{,}17 \times 10^{-14}} = 104
	\end{align}
	
	Das Verhältnis ist berechnet, kein Messwert; gemessen ist nur $\rho_{\text{CMB}}$. Theoretischer Ausdruck (102{,}8) und direkte Berechnung mit dem gemessenen $\rho_{\text{CMB}}$ (104) weichen um rund 1\% voneinander ab; das ist dieselbe Formel mit gerundetem $L_\xi$.
	
	\subsection{$\xi$-Feld als universelles Vakuum}
	
	\begin{important}
		Das $\xi$-Feld manifestiert sich sowohl in der freien CMB-Strahlung als auch im geometrisch beschränkten Casimir-Vakuum. In der Deutung dieses Dokuments ist das eine Skalenaussage \textbf{[S]}, kein Beleg für die physikalische Realität des $\xi$-Feldes.
	\end{important}
	
	Bei der charakteristischen $\xi$-Längenskala $L_\xi$ fallen CMB-Vakuum-Energiedichte und Casimir-Energiedichte nicht zusammen (Verhältnis rund $10^2$); Gleichheit liegt bei einem größeren Plattenabstand:
	
	\begin{align}
		\text{Freies Vakuum:} \quad &\rho_{\text{CMB}} \approx 2{,}0 \times 10^{-15}\ \text{eV}^4 \\
		\text{Beschränktes Vakuum:} \quad &|\rho_{\text{Casimir}}| = \frac{\pi^2}{720 d^4}
	\end{align}
	
	\section{Kosmische Rotverschiebung ohne Expansion}
	
	\subsection{$\xi$-Feld-Mechanismus}
	
	\begin{revolutionary}[These]
		\textbf{[S]} Nach dieser These entsteht die beobachtete kosmische Rotverschiebung nicht durch räumliche Expansion, sondern durch fraktale Wegverlängerung im omnipräsenten $\xi$-Feld.
	\end{revolutionary}
	
	Die Rotverschiebung ist achromatisch (K6): Die fraktale Wegverlängerung ist rein geometrisch und streckt alle Wellenlängen mit demselben Faktor, $1+z = e^{\xi x}$, unabhängig von $\lambda$; die Skala der Wegverlängerung ist über $H_0$ kalibriert (P39).
	
	Ein energieabhängiger Energieverlust der Photonen, etwa $dE/dx = -\xi E^2/E_\xi$, gilt nicht, weil er eine wellenlängenabhängige Rotverschiebung ergäbe; die beobachtete kosmologische Rotverschiebung ist achromatisch.
	
	\section{Strukturbildung im statischen $\xi$-Universum}
	
	\subsection{Kontinuierliche Strukturentwicklung}
	
	Im statischen T0-Universum \textbf{[S]} erfolgt Strukturbildung kontinuierlich ohne Urknall-Beschränkungen:
	
	\begin{equation}
		\frac{d\rho}{dt} = -\nabla \cdot (\rho \mathbf{v}) + S_\xi(\rho, T, \xi)
	\end{equation}
	
	wobei $S_\xi$ der $\xi$-Feld-Quellterm für kontinuierliche Materie/Energie-Transformation ist.
	
	\subsection{$\xi$-unterstützte kontinuierliche Schöpfung}
	
	Das $\xi$-Feld ermöglicht kontinuierliche Materie/Energie-Transformation:
	
	\begin{align}
		\text{Quantenvakuum} &\xrightarrow{\xi} \text{Virtuelle Teilchen} \\
		\text{Virtuelle Teilchen} &\xrightarrow{\xi^2} \text{Reale Teilchen} \\
		\text{Reale Teilchen} &\xrightarrow{\xi^3} \text{Atomkerne} \\
		\text{Atomkerne} &\xrightarrow{\text{Zeit}} \text{Sterne, Galaxien}
	\end{align}
	
	Die Energiebilanz wird aufrechterhalten durch:
	\begin{equation}
		\rho_{\text{gesamt}} = \rho_{\text{Materie}} + \rho_{\xi\text{-Feld}} = \text{konstant}
	\end{equation}
	
	\section{Dimensionslose $\xi$-Hierarchie}
	
	\subsection{Energieskalenverhältnisse}
	
	Die hier verwendeten $\xi$-Beziehungen reduzieren sich auf exakte mathematische Verhältnisse:
	
	\begin{longtable}{lcc}
		\caption{Dimensionslose $\xi$-Verhältnisse} \\
		\toprule
		\textbf{Verhältnis} & \textbf{Ausdruck} & \textbf{Wert} \\
		\midrule
		\endfirsthead
		\multicolumn{3}{c}{\tablename\ \thetable{} -- Fortsetzung} \\
		\toprule
		\textbf{Verhältnis} & \textbf{Ausdruck} & \textbf{Wert} \\
		\midrule
		\endhead
		Temperatur & $\frac{T_{\text{CMB}}}{E_\xi}$ & $3{,}13 \times 10^{-8}$ \\
		Relation (einheitenabhängig) & $\frac{16}{9}\xi^2$ & $3{,}16 \times 10^{-8}$ \\
		Länge & $\frac{\ell_{\xi}}{L_\xi}$ & $\xi^{-1/4}$ \\
		Casimir-CMB & $\frac{|\rho_{\text{Casimir}}|}{\rho_{\text{CMB}}}$ & $\frac{\pi^2 \times 10^4}{960}$ \\
		\bottomrule
	\end{longtable}
	
	\begin{important}
		Die hier verwendeten $\xi$-Beziehungen bestehen aus exakten mathematischen Verhältnissen:
		\begin{itemize}
			\item Brüche: $\frac{4}{3}$, $\frac{3}{4}$, $\frac{16}{9}$
			\item Zehnerpotenzen: $10^{-4}$, $10^3$, $10^4$
			\item Mathematische Konstanten: $\pi^2$
		\end{itemize}
		Willkürliche Dezimalzahlen treten in diesen Ausdrücken nicht auf; sie werden aus $\xi$ und ganzzahligen Strukturwahlen (motivierten Setzungen) gebildet.
	\end{important}
	
	\section{Experimentelle Vorhersagen und Tests}
	
	\subsection{Präzisionsmessungen der Gravitationskonstante}
	
	Die T0-Theorie sagt vorher:
	\begin{equation}
		G_{\text{T0}} = 6{,}6745 \times 10^{-11} \text{ m}^3/(\text{kg} \cdot \text{s}^2) \quad (+0{,}003\,\%\ \text{zu CODATA})
	\end{equation}
	
	Die Korpusformel $G = \xi^2/(4m_e)\cdot C_{\text{conv}}\cdot K_{\text{frak}}$ liefert $6{,}6745\times10^{-11}$. $C_{\text{conv}}$ ist die SI-Umrechnung der Ein-Anker-Kette (Dok.~383); ihr Zahlenwert $7{,}783\times10^{-3}$ enthält zusätzlich den kleinen Rest dieser Kette (Kettenwert $7{,}58\times10^{-3}$). Die Form von $G$ und seine Größenordnung folgen aus $\xi$ und dem einen Anker (numerisch geprüft in Dok.~383); mit dem Anker $m_e$ liegt $G$ bei $-2{,}6\,\%$, mit $v$ bei $-0{,}46\,\%$ (Dok.~383). Die Übereinstimmung auf $0{,}003\,\%$ ergibt sich erst mit diesem Rest und gilt nicht als eigener Präzisionsbeleg.
	
	\subsection{Casimir-Kraft-Anomalien}
	
	\begin{experiment}
		Vorhersage: Casimir-Kraft-Anomalien bei charakteristischer $\xi$-Längenskala
		\begin{itemize}
			\item Standard-Casimir-Gesetz: $F \propto d^{-4}$
			\item $\xi$-Feld-Modifikationen bei $d = L_\xi = 10^{-4}$ m
			\item Messbare Abweichungen durch $\xi$-Vakuum-Kopplung
		\end{itemize}
	\end{experiment}
	
	\subsection{Elektromagnetische Resonanz}
	
	Maximale $\xi$-Feld-Photon-Kopplung bei charakteristischer Frequenz:
	\begin{equation}
		\nu_\xi = \frac{c}{L_\xi} \approx 3 \times 10^{12} \text{ Hz} = 3 \text{ THz}
	\end{equation}

	Dabei ist $1/L_\xi = 10^4$~m$^{-1}$ eine Wellenzahl; die Frequenz ist $c/L_\xi$.
	
	Bei dieser Frequenz sollten elektromagnetische Anomalien auftreten.
	
	\section{Kosmologische Konsequenzen}
	
	\subsection{Kosmologische Probleme im T0-Bild}
	
	Die folgende Tabelle zeigt, wie sich bekannte Probleme der Standardkosmologie im statischen T0-Modell darstellen \textbf{[S]}. Die Einträge sind Deutungen im kosmischen Sektor, den der heutige Korpus ausklammert (P39), keine hergeleiteten Lösungen:
	
\begin{table}[ht]
	\centering
	\caption{Kosmologische Probleme: Standard vs. T0}
	\label{tab:cosmological_problems}
	\resizebox{\textwidth}{!}{%
		\begin{tabular}{lcc}
			\toprule
			\textbf{Problem} & \textbf{$\Lambda$CDM} & \textbf{T0-Deutung} \\
			\midrule
			Horizontproblem & Inflation erforderlich & Unendliche kausale Konnektivität \\
			Flachheitsproblem & Feinabstimmung & Geometrie stabilisiert über unendliche Zeit \\
			Monopolproblem & Topologische Defekte & Defekte dissipieren über unendliche Zeit \\
			Lithiumproblem & Nukleosynthese-Diskrepanz & Nukleosynthese über unbegrenzte Zeit \\
			Altersproblem & Objekte älter als Universum & Objekte können beliebig alt sein \\
			$H_0$-Spannung & 9\% Diskrepanz & Kein $H_0$ im statischen Universum \\
			Dunkle Energie & 69\% der Energiedichte & Nicht erforderlich \\
			\bottomrule
		\end{tabular}%
	}
\end{table}
	
	\subsection{Parameterreduktion}
	
	\begin{revolutionary}[Parameterzählung]
		Parameterzählung: Standardmodell und $\Lambda$CDM zusammen verwenden 25+ Parameter, die T0-Theorie den einzigen Parameter $\xi$ (dazu ganzzahlige Strukturwahlen als motivierte Setzungen; SI-Anker dienen nur der Umrechnung).
		\begin{itemize}
			\item Standardmodell der Teilchenphysik: 19+ Parameter
			\item $\Lambda$CDM-Kosmologie: 6 Parameter
			\item T0-Theorie: 1 Parameter ($\xi$)
		\end{itemize}
		Das entspricht rechnerisch einer Reduktion um 96\%. Der $\Lambda$CDM-Anteil betrifft den kosmischen Sektor, den der heutige Korpus ausklammert (P39).
	\end{revolutionary}
	
	\section{Schlussfolgerungen}

	\subsection{Das Vakuum ist das $\xi$-Feld}
	
	\begin{important}
		Zusammenfassung der T0-Deutung (kosmologische Punkte \textbf{[S]}):
		\begin{itemize}
			\item Das Vakuum ist identisch mit dem $\xi$-Feld
			\item Die CMB ist die Strahlung dieses Vakuums bei charakteristischer Temperatur
			\item Die Casimir-Kraft entsteht durch geometrische Beschränkung desselben Vakuums
			\item Gravitation folgt aus der $\xi$-Geometrie
			\item Kosmische Rotverschiebung entsteht achromatisch durch fraktale Wegverlängerung im $\xi$-Feld
		\end{itemize}
	\end{important}
	
	\subsection{Mathematische Struktur}
	
	Kennzeichen des hier vorgestellten Ansatzes:
	\begin{enumerate}
		\item \textbf{Universelle $\xi$-Skalierung}: Die behandelten Phänomene werden zurückgeführt auf $\xi = \frac{4}{3} \times 10^{-4}$
		\item \textbf{Statisches Modell} \textbf{[S]}: Kein Urknall, keine Expansion, ewige Existenz
		\item \textbf{Zeit-Energie-Konsistenz}: Respektiert fundamentale Quantenmechanik
		\item \textbf{Dimensionale Konsistenz}: Vollständig in natürlichen Einheiten formuliert
		\item \textbf{Einheitenunabhängige Physik}: Exakte mathematische Verhältnisse
	\end{enumerate}
	
	\begin{revolutionary}[Fazit]
		Die T0-Theorie bietet einen in natürlichen Einheiten formulierten Gegenentwurf zur expansionsbasierten Kosmologie und beschreibt die hier behandelten kosmischen Phänomene mit einer einzigen fundamentalen Konstante in einem statischen, ewig existierenden Universum \textbf{[S]}. Dieser kosmologische Teil bleibt spekulativ, solange der kosmische Sektor ausgeklammert ist (P39).
	\end{revolutionary}
	
	Die Übereinstimmungen zwischen theoretischen Ausdrücken und Messwerten sind unterschiedlich belastbar: Form und Größenordnung der Gravitationskonstante folgen aus $\xi$ und einem Anker, die Übereinstimmung auf $0{,}003\,\%$ ist kein eigener Präzisionsbeleg; die CMB-Temperatur wird einheitenabhängig auf etwa 1\,\% getroffen; das Casimir-CMB-Verhältnis ist eine Identität. Die kosmologische Deutung bleibt spekulativ.
	

	
	\begin{thebibliography}{20}
		
		\bibitem{t0_lagrangian_de}
		Pascher, Johann (2025). 
		\textit{Vereinfachte Lagrange-Dichte und Zeit-Massen-Dualität in der T0-Theorie}. 
		T0-Theorie Projekt. 
		\url{https://jpascher.github.io/T0-Time-Mass-Duality/2/pdf/129_lagrandian-einfach_De.pdf}
		
		\bibitem{t0_lagrangian_en}
		Pascher, Johann (2025). 
		\textit{Simplified Lagrangian Density and Time-Mass Duality in T0-Theory}. 
		T0-Theory Project. 
		\url{https://jpascher.github.io/T0-Time-Mass-Duality/2/pdf/129_lagrandian-einfach_En.pdf}
		
		\bibitem{t0_cosmos_de}
		Pascher, Johann (2025). 
		\textit{T0-Modell: Ein vereinheitlichtes, statisches, zyklisches, dunkle-Materie-freies und dunkle-Energie-freies Universum}. 
		T0-Theorie Projekt. 
		\url{https://jpascher.github.io/T0-Time-Mass-Duality/2/pdf/063_cosmic_De.pdf}
		
		\bibitem{t0_cosmos_en}
		Pascher, Johann (2025). 
		\textit{T0-Model: A unified, static, cyclic, dark-matter-free and dark-energy-free universe}. 
		T0-Theory Project. 
		\url{https://jpascher.github.io/T0-Time-Mass-Duality/2/pdf/063_cosmic_En.pdf}
		
		\bibitem{t0_cmb_de}
		Pascher, Johann (2025). 
		\textit{Temperatureinheiten in natürlichen Einheiten: T0-Theorie und statisches Universum}. 
		T0-Theorie Projekt. 
		\url{https://jpascher.github.io/T0-Time-Mass-Duality/2/pdf/061_TempEinheitenCMB_De.pdf}
		
		\bibitem{t0_cmb_en}
		Pascher, Johann (2025). 
		\textit{Temperature Units in Natural Units: T0-Theory and Static Universe}. 
		T0-Theory Project. 
		\url{https://jpascher.github.io/T0-Time-Mass-Duality/2/pdf/061_TempEinheitenCMB_En.pdf}
		
		\bibitem{t0_gravitation_en}
		Pascher, Johann (2025). 
		\textit{Geometric Determination of the Gravitational Constant: From the T0-Model}. 
		T0-Theorie Projekt. 
		\url{https://jpascher.github.io/T0-Time-Mass-Duality/2/pdf/127_gravitationskonstnte_En.pdf}
		
		\bibitem{t0_redshift_de}
		Pascher, Johann (2025). 
		\textit{T0-Theorie: Wellenlängenabhängige Rotverschiebung ohne Distanzannahmen}. 
		T0-Theorie Projekt. 
		\url{https://jpascher.github.io/T0-Time-Mass-Duality/2/pdf/065_redshift_deflection_De.pdf}
		
		\bibitem{t0_redshift_en}
		Pascher, Johann (2025). 
		\textit{T0-Theory: Wavelength-Dependent Redshift without Distance Assumptions}. 
		T0-Theorie Projekt. 
		\url{https://jpascher.github.io/T0-Time-Mass-Duality/2/pdf/065_redshift_deflection_En.pdf}
		
		\bibitem{heisenberg1927}
		Heisenberg, W. (1927). 
		\textit{Über den anschaulichen Inhalt der quantentheoretischen Kinematik und Mechanik}. 
		Zeitschrift für Physik, 43(3-4), 172--198.
		
		\bibitem{planck2020}
		Planck Collaboration (2020). 
		\textit{Planck 2018 results. VI. Cosmological parameters}. 
		Astronomy \& Astrophysics, 641, A6. 
		\url{https://doi.org/10.1051/0004-6361/201833910}
		
		\bibitem{codata2018}
		CODATA (2018). 
		\textit{The 2018 CODATA Recommended Values of the Fundamental Physical Constants}. 
		National Institute of Standards and Technology. 
		\url{https://physics.nist.gov/cuu/Constants/}
		
		\bibitem{casimir1948}
		Casimir, H. B. G. (1948). 
		\textit{On the attraction between two perfectly conducting plates}. 
		Proceedings of the Royal Netherlands Academy of Arts and Sciences, 51(7), 793--795.
		
		\bibitem{muon_g2_2021}
		Muon g-2 Collaboration (2021). 
		\textit{Measurement of the Positive Muon Anomalous Magnetic Moment to 0.46 ppm}. 
		Physical Review Letters, 126(14), 141801. 
		\url{https://doi.org/10.1103/PhysRevLett.126.141801}
		
		\bibitem{riess2022}
		Riess, A. G., et al. (2022). 
		\textit{A Comprehensive Measurement of the Local Value of the Hubble Constant with 1 km s$^{-1}$ Mpc$^{-1}$ Uncertainty from the Hubble Space Telescope and the SH0ES Team}. 
		The Astrophysical Journal Letters, 934(1), L7. 
		\url{https://doi.org/10.3847/2041-8213/ac5c5b}
		
		\bibitem{jwst_early}
		Naidu, R. P., et al. (2022). 
		\textit{Two Remarkably Luminous Galaxy Candidates at z $\approx$ 11--13 Revealed by JWST}. 
		The Astrophysical Journal Letters, 940(1), L14. 
		\url{https://doi.org/10.3847/2041-8213/ac9b22}
		
		\bibitem{cobe1992}
		COBE Collaboration (1992). 
		\textit{Structure in the COBE differential microwave radiometer first-year maps}. 
		The Astrophysical Journal Letters, 396, L1--L5. 
		\url{https://doi.org/10.1086/186504}
		
		\bibitem{sparnaay1958}
		Sparnaay, M. J. (1958). 
		\textit{Measurements of attractive forces between flat plates}. 
		Physica, 24(6-10), 751--764. 
		\url{https://doi.org/10.1016/S0031-8914(58)80090-7}
		
		\bibitem{lamoreaux1997}
		Lamoreaux, S. K. (1997). 
		\textit{Demonstration of the Casimir force in the 0.6 to 6 $\mu$m range}. 
		Physical Review Letters, 78(1), 5--8. 
		\url{https://doi.org/10.1103/PhysRevLett.78.5}
		
		\bibitem{einstein1915}
		Einstein, A. (1915). 
		\textit{Die Feldgleichungen der Gravitation}. 
		Sitzungsberichte der Preußischen Akademie der Wissenschaften, 844--847.
		
	\end{thebibliography}
	